\documentclass[11pt]{article}
\usepackage[a4paper,margin=24mm]{geometry}
\usepackage{fontspec}
\usepackage{xcolor}
\usepackage{amsmath,amssymb}
\usepackage{booktabs,longtable,array}
\usepackage[hidelinks]{hyperref}
\setmainfont[Path=fonts/,Extension=.otf,UprightFont=*-Regular,BoldFont=*-Bold]{NotoSerifCJKkr}
\definecolor{reviewercolor}{RGB}{38,82,126}
\definecolor{responsecolor}{RGB}{36,113,78}
\definecolor{revisioncolor}{RGB}{145,82,24}
\setlength{\parindent}{0pt}
\setlength{\parskip}{0.55em}
\setlength{\fboxsep}{7pt}
\setlength{\emergencystretch}{3em}
\renewcommand{\familydefault}{\rmdefault}
\raggedright
\newcommand{\fieldheading}[2]{\par\medskip\noindent{\color{#1}\rule{3pt}{1.25em}}\hspace{0.55em}\textbf{#2}\par\smallskip}
\newenvironment{reviewitem}[1]{\subsection*{#1}}{\par\medskip\hrule\bigskip}
\newenvironment{reviewercomment}{\fieldheading{reviewercolor}{1. 리뷰어 코멘트}\begin{quote}}{\end{quote}}
\newenvironment{authorresponse}{\fieldheading{responsecolor}{2. 저자 답변}}{\par}
\newenvironment{manuscriptrevision}[1]{\fieldheading{revisioncolor}{3. 논문 반영 결과}\textbf{수정 위치:} #1\par\smallskip}{\par}
\newcommand{\manuscripttitle}{연속 인과관계(CoC) 주석의 시간·조건 일관성 검사 프레임워크: 적용 가능성과 근거 추적}
\newcommand{\manuscriptid}{확인 필요 (제공 기록에서 확인되지 않음)}
\newcommand{\revisionround}{확인 필요 (개정 소스 버전: revision-v01)}
\newcommand{\responsedate}{2026년 10월 12일}
% Generated from current manuscript auxiliary labels.
\expandafter\def\csname mloc@sec:intro\endcsname{1절(1쪽)}
\expandafter\def\csname mloc@fig:overview-pipeline\endcsname{그림 1(2쪽)}
\expandafter\def\csname mloc@sec:related\endcsname{2절(2쪽)}
\expandafter\def\csname mloc@sec:related-reasoning\endcsname{2.1절(2쪽)}
\expandafter\def\csname mloc@sec:related-rules\endcsname{2.2절(2쪽)}
\expandafter\def\csname mloc@sec:related-nl\endcsname{2.3절(2쪽)}
\expandafter\def\csname mloc@sec:related-rv\endcsname{2.4절(3쪽)}
\expandafter\def\csname mloc@sec:related-position\endcsname{2.5절(3쪽)}
\expandafter\def\csname mloc@sec:data\endcsname{3절(3쪽)}
\expandafter\def\csname mloc@sec:data-overview\endcsname{3.1절(3쪽)}
\expandafter\def\csname mloc@sec:provenance\endcsname{3.2절(3쪽)}
\expandafter\def\csname mloc@sec:review-protocol\endcsname{3.3절(3쪽)}
\expandafter\def\csname mloc@sec:problem\endcsname{4절(3쪽)}
\expandafter\def\csname mloc@tab:related-position\endcsname{표 1(4쪽)}
\expandafter\def\csname mloc@tab:evaluation-data\endcsname{표 2(4쪽)}
\expandafter\def\csname mloc@tab:provenance\endcsname{표 3(4쪽)}
\expandafter\def\csname mloc@sec:notation\endcsname{4.1절(4쪽)}
\expandafter\def\csname mloc@sec:error-types\endcsname{4.2절(4쪽)}
\expandafter\def\csname mloc@eq:hold-conflict\endcsname{식 2(4쪽)}
\expandafter\def\csname mloc@eq:missing-evidence\endcsname{식 6(4쪽)}
\expandafter\def\csname mloc@tab:model-scope\endcsname{표 4(5쪽)}
\expandafter\def\csname mloc@sec:semantics\endcsname{4.3절(5쪽)}
\expandafter\def\csname mloc@eq:active-update\endcsname{식 9(5쪽)}
\expandafter\def\csname mloc@sec:boundary-semantics\endcsname{4.4절(5쪽)}
\expandafter\def\csname mloc@tab:boundary-policy\endcsname{표 5(5쪽)}
\expandafter\def\csname mloc@sec:issue-definition\endcsname{4.5절(5쪽)}
\expandafter\def\csname mloc@tab:parser-rules\endcsname{표 6(6쪽)}
\expandafter\def\csname mloc@eq:response-contract\endcsname{식 13(6쪽)}
\expandafter\def\csname mloc@sec:method\endcsname{5절(6쪽)}
\expandafter\def\csname mloc@sec:workflow\endcsname{5.1절(6쪽)}
\expandafter\def\csname mloc@sec:extraction\endcsname{5.2절(6쪽)}
\expandafter\def\csname mloc@sec:models\endcsname{5.3절(6쪽)}
\expandafter\def\csname mloc@sec:queries\endcsname{5.3절(6쪽)}
\expandafter\def\csname mloc@sec:trace\endcsname{5.3절(6쪽)}
\expandafter\def\csname mloc@lst:condition-step\endcsname{목록 1(6쪽)}
\expandafter\def\csname mloc@fig:method-pipeline\endcsname{그림 2(7쪽)}
\expandafter\def\csname mloc@tab:worked-example\endcsname{표 7(7쪽)}
\expandafter\def\csname mloc@tab:model-variables\endcsname{표 8(7쪽)}
\expandafter\def\csname mloc@fig:condition-observer\endcsname{그림 3(7쪽)}
\expandafter\def\csname mloc@lst:queries\endcsname{목록 2(7쪽)}
\expandafter\def\csname mloc@tab:query-meaning\endcsname{표 9(7쪽)}
\expandafter\def\csname mloc@sec:boundary-implementation\endcsname{5.4절(7쪽)}
\expandafter\def\csname mloc@sec:model-scope\endcsname{5.5절(8쪽)}
\expandafter\def\csname mloc@sec:experiments\endcsname{6절(8쪽)}
\expandafter\def\csname mloc@sec:legacy-test\endcsname{6.1절(8쪽)}
\expandafter\def\csname mloc@sec:scenario-design\endcsname{6.2절(8쪽)}
\expandafter\def\csname mloc@sec:retest-design\endcsname{6.3절(8쪽)}
\expandafter\def\csname mloc@sec:boundary-design\endcsname{6.4절(8쪽)}
\expandafter\def\csname mloc@sec:bounded-design\endcsname{6.5절(8쪽)}
\expandafter\def\csname mloc@tab:comparison-paths\endcsname{표 10(9쪽)}
\expandafter\def\csname mloc@tab:scenario-variants\endcsname{표 11(9쪽)}
\expandafter\def\csname mloc@sec:human-design\endcsname{6.6절(9쪽)}
\expandafter\def\csname mloc@sec:aux-design\endcsname{6.7절(9쪽)}
\expandafter\def\csname mloc@tab:regression\endcsname{표 12(9쪽)}
\expandafter\def\csname mloc@sec:results\endcsname{7절(9쪽)}
\expandafter\def\csname mloc@sec:regression-results\endcsname{7.1절(9쪽)}
\expandafter\def\csname mloc@sec:first-results\endcsname{7.2절(9쪽)}
\expandafter\def\csname mloc@tab:first-results\endcsname{표 13(10쪽)}
\expandafter\def\csname mloc@tab:failures\endcsname{표 14(10쪽)}
\expandafter\def\csname mloc@sec:retest-results\endcsname{7.3절(10쪽)}
\expandafter\def\csname mloc@sec:agreement-results\endcsname{7.4절(10쪽)}
\expandafter\def\csname mloc@tab:retest\endcsname{표 15(10쪽)}
\expandafter\def\csname mloc@tab:implementation\endcsname{표 16(10쪽)}
\expandafter\def\csname mloc@tab:boundary-results\endcsname{표 17(10쪽)}
\expandafter\def\csname mloc@tab:bounded-results\endcsname{표 18(10쪽)}
\expandafter\def\csname mloc@sec:boundary-results\endcsname{7.5절(10쪽)}
\expandafter\def\csname mloc@sec:bounded-results\endcsname{7.6절(10쪽)}
\expandafter\def\csname mloc@tab:agreement\endcsname{표 19(11쪽)}
\expandafter\def\csname mloc@tab:reviewer-distribution\endcsname{표 20(11쪽)}
\expandafter\def\csname mloc@sec:human-results\endcsname{7.7절(11쪽)}
\expandafter\def\csname mloc@sec:trajectory-results\endcsname{7.8절(11쪽)}
\expandafter\def\csname mloc@sec:discussion\endcsname{8절(11쪽)}
\expandafter\def\csname mloc@sec:discussion-results\endcsname{8.1절(11쪽)}
\expandafter\def\csname mloc@sec:review-use\endcsname{8.2절(11쪽)}
\expandafter\def\csname mloc@tab:disagreement-examples\endcsname{표 21(12쪽)}
\expandafter\def\csname mloc@tab:trajectory-cases\endcsname{표 22(12쪽)}
\expandafter\def\csname mloc@tab:review-actions\endcsname{표 23(12쪽)}
\expandafter\def\csname mloc@sec:threats\endcsname{8.3절(12쪽)}
\expandafter\def\csname mloc@sec:future-work\endcsname{8.4절(12쪽)}
\expandafter\def\csname mloc@sec:conclusion\endcsname{9절(12쪽)}
\expandafter\def\csname mloc@sec:appendix-trajectory\endcsname{9절(12쪽)}
\expandafter\def\csname mloc@tab:limitations\endcsname{표 24(13쪽)}
\expandafter\def\csname mloc@tab:trajectory-details\endcsname{표 25(13쪽)}

% Generated from exact excerpts in the current main.pdf.
\expandafter\def\csname quotepages@R1-m3\endcsname{12쪽}
\expandafter\def\csname quotepages@R1-m2\endcsname{3쪽}
\expandafter\def\csname quotepages@R1-m1\endcsname{3쪽}
\expandafter\def\csname quotepages@R1-M6\endcsname{3쪽}
\expandafter\def\csname quotepages@R1-M5\endcsname{11쪽}
\expandafter\def\csname quotepages@R1-M4\endcsname{10쪽}
\expandafter\def\csname quotepages@R1-M3\endcsname{11쪽}
\expandafter\def\csname quotepages@R1-M2\endcsname{6쪽}
\expandafter\def\csname quotepages@R1-M1\endcsname{12쪽}
\expandafter\def\csname quotepages@R2-7\endcsname{2쪽}
\expandafter\def\csname quotepages@R2-6\endcsname{3쪽}
\expandafter\def\csname quotepages@R2-5\endcsname{1쪽}
\expandafter\def\csname quotepages@R2-4\endcsname{2쪽}
\expandafter\def\csname quotepages@R2-3\endcsname{1쪽}
\expandafter\def\csname quotepages@R2-2\endcsname{1쪽}
\expandafter\def\csname quotepages@R2-1.2\endcsname{1쪽}
\expandafter\def\csname quotepages@R2-1.1\endcsname{1쪽}
\expandafter\def\csname quotepages@R3-11\endcsname{11쪽}
\expandafter\def\csname quotepages@R3-10\endcsname{9쪽}
\expandafter\def\csname quotepages@R3-9\endcsname{11쪽}
\expandafter\def\csname quotepages@R3-8\endcsname{10쪽}
\expandafter\def\csname quotepages@R3-7\endcsname{10쪽}
\expandafter\def\csname quotepages@R3-6\endcsname{6쪽}
\expandafter\def\csname quotepages@R3-5\endcsname{6쪽}
\expandafter\def\csname quotepages@R3-4\endcsname{5쪽}
\expandafter\def\csname quotepages@R3-3\endcsname{4쪽}
\expandafter\def\csname quotepages@R3-2\endcsname{3쪽}
\expandafter\def\csname quotepages@R3-1\endcsname{3쪽}

\newcommand{\mloc}[1]{\csname mloc@#1\endcsname}
\newcommand{\quotepages}[1]{\csname quotepages@#1\endcsname}
\newcommand{\responsefront}[2]{%
\begin{center}{\LARGE\bfseries 심사 의견 답변서\par}\vspace{1em}{\large\bfseries\manuscripttitle\par}\end{center}
\vspace{0.8em}
\begin{tabular}{@{}ll@{}}논문 접수번호: & \manuscriptid \\ 수정 차수: & \revisionround \\ 답변일: & \responsedate\end{tabular}
\section*{답변에 앞서}
편집위원장님과 심사위원님께,

연속 CoC의 검사 범위, 의미 변환과 평가 근거를 구체화할 수 있도록 의견을 주셔서 감사합니다.
이번 개정은 자연 오류 성능 대신 시간·조건 관계의 검사 적용 가능성과 근거 추적으로 주장을 제한했습니다.
기존 통제 시나리오의 최초 결과와 개발 후 동일 사례 재시험을 구분하고, 81건의 장면 경계 시험과 기존 사람 라벨의 AC1 재분석을 본문에 추가했습니다.
선별 없는 유한 직적 전수 시험은 고정 XML의 step 코드와 별도 수식 판정기를 대조한 구현 적합성 근거입니다.
두 경로에 같은 작성자가 관여했으므로 외부 독립 시험이나 자연 의미 정답 검증을 완료했다고 주장하지 않습니다. 사람 재평가도 수행하지 않았습니다.
개정 원고의 파란색은 원 저장소의 최초 제출본 대비 이번에 추가·수정한 내용을 나타냅니다.
검정색 clean 원고는 동일 내용입니다. 아래 절·쪽·표 번호는 이번 개정 원고를 기준으로 합니다.
접수번호와 실제 심사 수정 차수는 기록이 없어 확인 필요로 남겼습니다.

\section*{심사위원 #1}
\textbf{총평(요지)}\par
\begin{quote}#2\end{quote}}


\begin{document}

\responsefront{3}{연속 CoC의 상태 검사 접근은 의미가 있으나 사전 확인한 합성자료, 낮은 자연 검토 일치도와 자연 모순 부재가 실증을 제한한다. 의미 추출과 구현 설명, 관련 연구, 독립 시험과 평가 근거를 보완할 필요가 있다.}

\clearpage

\begin{reviewitem}{R3-1. 시간 논리와 실행 중 검증 연구}

\begin{reviewercomment}

2장 관련 연구: 시간 논리 기반 모니터링, runtime verification 및 기존 자율주행 설명 검증 연구를 보완하고, 제안 방법과 기존 연구의 기술적 차이를 구체적으로 제시해 주십시오.

\end{reviewercomment}

\begin{authorresponse}

시간 논리·실행 중 검증 문헌과 자율주행 설명·계획 연구를 1차 자료로 확인하고 비교표를 추가했습니다. 설명 충실성 연구(Is VLA Reasoning Faithful?)와 VLADriveBench의 관측·개입 검증도 저장 주석의 조건 연결 검사와 비교했습니다. 관측 경로 검사, LTL/TLTL 유한 접두 경로 판정, UPPAAL 상태 질의와 본 논문의 조건별 CoC 근거 추적을 입력·표현·출력으로 비교했습니다. 상태 기억·세 값·최초 위반은 기존 기술임을 인정했습니다. 또한 교통 규칙의 시간 논리 형식화와 규칙 검증(Esterle 등, Maierhofer 등, Wan 등), 자연어–시간 논리 변환(NL2TL, nl2spec), 추론 충실성과 VLA 추론–궤적 일관성 감시(Turpin 등, Lanham 등, Drive the Thoughts)를 2장에 추가하고 비교표에 반영했습니다. 이들은 사람이 정한 규칙, 개별 요구 문장 또는 한 장면의 설명–행동 정렬을 다루며, 본 논문은 자동 생성된 연속 주석에서 의무와 해제 근거를 연결한다는 점을 기술적 차이로 제시했습니다.

두 모델의 근거 부족은 해제 증거 또는 연결 근거의 부족이며 LTL의 미래 연장에 대한 미결정과 같다고 하지 않습니다. 도구 사용과 응용 대상의 차이를 설명하고, 직접 비교하지 않은 연구에 대한 우월성이나 최초 주장은 제거했습니다.

\end{authorresponse}

\begin{manuscriptrevision}{\mloc{sec:related-rv}; 발췌: \quotepages{R3-1}}
본문의 실제 수정 문장은 다음과 같습니다.

\begin{quote}
실행 중 검증(runtime verification)은 관측된 실행 경로가 주어진 명세를 만족하는지 검사한다~[15]. LTL/TLTL 모니터링은 유한 접두 경로에서 참ㆍ거짓ㆍ미결정을 구분하고 최초로 확정되는 시점을 다룬다~[16]. 따라서 이전 상태의 기억, 세 값의 판정과 최초 위반 추적 자체는 기존 기술이다.
\end{quote}
\end{manuscriptrevision}

\end{reviewitem}

\clearpage

\begin{reviewitem}{R3-2. 자료 버전·생성과 선정 절차}

\begin{reviewercomment}

3장 데이터와 증거 범위: CoC-Nusc 자료의 버전, CoC 생성 모델과 프롬프트, 품질 점수 산정 방법 및 403개 사건 중 290개와 134개를 선정한 기준을 명시해 제시. 또한 자연 검토 대상 27개 장면의 선정 절차를 구체적으로 설명.

\end{reviewercomment}

\begin{authorresponse}

3장을 평가 자료 개요, 자연 자료의 출처·선정, 사람 검토 절차·기록의 한계로 재구성했습니다. 설계자가 구성한 144건·81건과 무선별 전수 사건열을 핵심 평가 자료로 먼저 제시하고, 자연 자료가 쓰이는 사람 검토와 보조 반응 분석을 구분했습니다.

보관된 주석 중 자차 궤적 파일과 대응되는 94장면·403사건을 사용한 이유와, 309인접쌍이 사람 검토 후보 선정용이라는 사실을 명시했습니다. 403→290은 품질 통과, 290→134는 방향이 불명확하거나 복합인 156건의 제외이며 134=130탐지 대상+4이미 만족입니다. 반응 결과 125+4+5는 결과 절에만 두었습니다.

자료 설명서의 모델 계열과 실행별 생성 정보 미복원을 구분했습니다. 자연 검토는 보존된 원 선정 절차의 15후보장면·12비교장면이며 재실행으로 기존 검토 ID가 재현됩니다. 선정의 비무작위성과 생성 정보 결손 때문에 완전 재생성이나 대표성은 주장하지 않습니다.


\end{authorresponse}

\begin{manuscriptrevision}{\mloc{sec:provenance}; 발췌: \quotepages{R3-2}}
본문의 실제 수정 문장은 다음과 같습니다.

\begin{quote}
보조 반응 분석에서는 403사건 중 품질 표시를 통과한 290건을 먼저 남겼다. 이 중 복합 행동이나 방향이 불명확한 156건을 제외한 134건을 사용했다.
\end{quote}
\end{manuscriptrevision}

\end{reviewitem}

\clearpage

\begin{reviewitem}{R3-3. P4 검사기 처리 오류 분리}

\begin{reviewercomment}

4장 문제 정의의 P1–P4: P1–P3은 주석 내용의 불일치인 반면, P4는 검사기의 상태 갱신 오류에 해당됨. P4를 주석의 일관성 오류와 분리하여 별도의 ‘검사기 처리 오류’로 정의하고 관련 수식과 설명을 수정.

\end{reviewercomment}

\begin{authorresponse}

지적을 수용했습니다. P1–P3은 입력 행동·요구 문제, P4는 해제 확인 뒤 이전 활성 의무를 상태에 남기는 처리 오류 $F_4$로 구분했습니다. 기존 40쌍도 P1–P3 30쌍과 P4 10처리 시험으로 분리했습니다. 기존 전체 이력의 40/40을 단일한 자연 주석 오류 검출 정확도로 쓰지 않습니다.

문제별 구현 범위 표를 추가했습니다. 조건 연결 모델과 경계 확장 모델에는 별도 해제 명령 입력이 없어 P2 자체를 평가하지 않습니다. P1의 진행 충돌과 P2의 조기 해제는 동치가 아닙니다. 기존 P2의 대용 규칙과 P4를 일으키는 반복·해제 조건 관측 입력은 7.1절로 옮겼고, 관측값을 해제 명령으로 설명하지 않았습니다. 144건에는 P4 고장 주입이 없으므로 그 성공을 모든 P1–P4 지원 성능으로 소급하지 않습니다.

% Supplemental evidence 2026-10-12
전수 적합성 시험에서도 P4는 참 증거로 해제된 활성 의무를 잘못 남기는 처리 고장으로 별도 주입했습니다. 길이 1–2의 5,256건 중 상태 불일치 1,884건을 보존했고 주 시험 분모에 합산하지 않았습니다. 입력 생애주기의 재활성화와 P4 코드 고장을 구분했습니다.
% End supplemental evidence

\end{authorresponse}

\begin{manuscriptrevision}{\mloc{sec:error-types}; 발췌: \quotepages{R3-3}}
본문의 실제 수정 문장은 다음과 같습니다.

\begin{quote}
P1--P3은 입력의 행동ㆍ요구 관계에 관한 문제이고 P4는 검사기의 상태 처리 문제다. P5는 별도의 속도 반응 시간 검사다. 기존 40쌍을 이 구분으로 재집계한다.
\end{quote}
\end{manuscriptrevision}

\end{reviewitem}

\clearpage

\begin{reviewitem}{R3-4. 미상·장면 종료·복수 의무}

\begin{reviewercomment}

4장 및 5장의 Unknown 처리: 종료 조건이 Unknown일 때 요구가 언제까지 유지되는지, 장면 종료 시 어떻게 처리되는지, 복수 요구가 동시에 활성화된 경우 어떤 순서로 판정하는지를 상태전이표나 의사코드로 제시.

\end{reviewercomment}

\begin{authorresponse}

조건 연결 모델의 기호·초기 상태·갱신식과 경계 확장 모델의 정책을 구분해 4장을 다시 구성했습니다. 조건 연결 모델은 미상 입력으로 알려진 증거를 덮지 않고 의무를 유지하며, 참 증거가 해당 의무만 지웁니다. 대기 중 미상 자체는 문제로 누적하지 않지만 해제 미상 상태에서 GO하면 근거 부족입니다. 저장된 선행조건이 거짓이면 확정 모순이며 미상 자체는 별도의 근거 부족 사유로 세지 않습니다.

조건 연결 모델은 선행조건 위반을 먼저, 의무 등록 순서대로 검사합니다. 확정 모순과 근거 부족을 각각 기록하되 최종 판정은 확정 모순을 우선하고 대표 최초 의무는 한 개를 반환합니다. 장면 끝에서 미래 해제를 못 봤다는 이유만으로 모순을 추가하지 않으며 장면마다 상태를 초기화합니다.

증거 재사용 정책도 명시했습니다. 조건 연결 모델은 의무 생성 때 이전 증거를 초기화하지 않습니다. 144건에는 같은 ID의 재생성이 없지만 전수시험에는 포함되며, 불일치 0은 이 정책의 구현 적합성입니다. 경계 확장 모델은 생성·재활성화 때 증거를 미상으로 초기화하고, 연결 미확인 동안 관측을 반영하지 않으며, 범위 종료 시 의무와 증거를 정리합니다. 이후 VERIFIED나 새 활성화로 연결이 확인되면 새 관측을 사용할 수 있습니다. 경계 확장 모델은 선행조건을 검사하지 않으므로, 확인된 지속 연결만으로 두 모델이 동등해지지는 않습니다. 실제 객체 추적·증거 유효 구간 추정은 두 모델 모두 수행하지 않습니다.

% Supplemental evidence 2026-10-12
경계 확장 모델의 VERIFIED/ENDED/UNCERTAIN, 순간 관측 만료, 지속 구간 근거, 재활성화 초기화를 본문에 추가했습니다. 24건에서 확정 모순과 근거 부족이 함께 발생했고 두 기록을 보존하며 최종 판정은 확정 모순을 우선했습니다. 대상·범위는 제공된 입력이며 영상에서 자동 추적한 결과가 아닙니다.
% End supplemental evidence

\end{authorresponse}

\begin{manuscriptrevision}{\mloc{sec:boundary-semantics}; 발췌: \quotepages{R3-4}}
본문의 실제 수정 문장은 다음과 같습니다.

\begin{quote}
갱신 순서는 경계 정책, 새 의무 활성화, 대응 증거 갱신과 해제, 진행 판정이다. 같은 사건의 새 의무는 새 적용 관계를 시작하며 이전 해제값을 재사용하지 않는다. 알려진 미해제 상태의 진행은 확정 모순, 연결 미확인 또는 해제 미상의 진행은 근거 부족이다.
\end{quote}
\end{manuscriptrevision}

\end{reviewitem}

\clearpage

\begin{reviewitem}{R3-5. 추출 주체와 전체 변환 사례}

\begin{reviewercomment}

5장 제안 방법: CoC 문장에서 행동, 지속 요구, 종료 조건, 선행조건 및 True/False/Unknown 증거를 추출하는 주체와 방법이 불분명. 수작업, 규칙 기반 또는 LLM 기반 여부와 세부 알고리즘을 명시하고, 실제 CoC 문장이 검사 입력으로 변환되는 전체 사례를 추가.

\end{reviewercomment}

\begin{authorresponse}

기존 자동 경로는 제한된 영어 정규식 파서이며 LLM 기반이 아님을 명시했습니다. 행동 동의어·명령 문맥·조건 구절·단계 연결·대상 누락·모호 처리 규칙을 표로 정리했습니다. 144건 통제 시나리오의 주평가에는 설계자가 원문과 함께 작성한 프레임을 사용했으며, 자동 텍스트 경로는 144건 전부 미상 출력임을 입력 설명 첫머리에도 밝혔습니다.

작성 보행자 원문의 세 변형을 조건별 프레임, 사례 전체 관측 배열, 다섯 질의의 결과와 확정 모순/근거 부족/양립까지 연결했습니다. 관측 배열은 공통 e0과 각 e1을 포함한 전체 두 사건의 값입니다. A와 B를 서로 다른 조건에 연결하고, 최초 문제 사건·의무에서 관련 원문 구절과 연결하여 제시하는 과정을 설명했습니다. 이 사례는 설계자 작성 자료이며 자연 자료에서 발견한 모순이나 자동 NLP의 성공 사례로 제시하지 않습니다.

% Supplemental evidence 2026-10-12
자연 CoC 90장면·399원문의 자동 파서 집계와 scene-0051의 실제 처리 사례를 추가했습니다. 265개 행동 단계 중 구문 해석 완료는 219개이며, 해제 조건을 찾거나 조건의 참·거짓을 확정한 성능으로 해석하지 않습니다. scene-0051 문장 하나의 기존 검사 출력은 양립이지만 후속 진행이 없어 충돌을 검사하지 않은 경우임을 함께 밝혔습니다. 자연 문장 자동 처리와 설계자 작성 프레임 사례를 같은 표의 (a)/(b)로 나눴습니다.
% End supplemental evidence

\end{authorresponse}

\begin{manuscriptrevision}{\mloc{sec:extraction}; 발췌: \quotepages{R3-5}}
본문의 실제 수정 문장은 다음과 같습니다.

\begin{quote}
True/False는 작성 문장이 명시한 조건 상태이며 사람의 독립 센서 정답이 아니다.
\end{quote}
\end{manuscriptrevision}

\end{reviewitem}

\clearpage

\begin{reviewitem}{R3-6. Python–UPPAAL 상태·속성과 일치성}

\begin{reviewercomment}

5장 Python–UPPAAL 구현: Python 검사기와 UPPAAL 모델의 상태, 전이 조건, 검증 속성 및 질의식을 구체적으로 제시. 두 구현의 결과 비교가 규칙의 타당성 검증이 아니라 구현 일치성 확인이라는 점도 더욱 명확히 기술.

\end{reviewercomment}

\begin{authorresponse}

조건 연결 모델의 입력·상태·기록 변수를 표로 정리하고, 실제 XML 상태 갱신 함수의 등록·증거 갱신·해제·진행 검사 순서를 의사코드로 제시했습니다. committed Run의 자기 전이와 Done 종료 조건, 실제 다섯 질의와 각 질의의 의미를 본문에 남겼습니다. 의무 상태와 확정 모순ㆍ근거 부족 기록은 위치가 아닌 변수이며, 고정된 입력마다 이산 상태 갱신 경로는 하나입니다. 시계와 환경 프로세스가 없어 환경의 여러 행동이나 정량적 시간을 검증한 결과로 해석하지 않습니다.

기존 단일 의무 검사기의 6상태와 대기 미상 누적 정책은 7.3절의 실패 분석으로 옮기고, 5.5절에는 비교 경로별 표현과 정책을 요약했습니다. 기존 Python--UPPAAL의 판정·오류 종류·최초 위치는 144/144 일치하지만 미상 사유는 최초 반례에서 24건이 달랐습니다. 종료 질의로 뒤의 누적 사유까지 읽으면 144/144가 일치하며 원 파서는 변경하지 않았습니다. 이는 같은 규칙의 구현 일치성이며 규칙·자연어 해석 자체의 독립 타당성 증명은 아닙니다.

% Supplemental evidence 2026-10-12
신규 81건의 배열·경계 갱신과 종료 시점의 확정 모순ㆍ근거 부족 질의를 기존 144건 모델과 구분했습니다. 추가 전수 시험은 고정 XML의 상태 갱신 코드를 C++로 실행하고 전체 접두열에서 계산하는 별도 수식 판정과 상태·누적 문제·최초 사건/의무를 대조한 것입니다. 35,042,928건 불일치 0은 유한 범위의 구현 적합성이지 규칙의 현실 타당성이나 UPPAAL 탐색 엔진의 새 검증이 아닙니다.
% End supplemental evidence

\end{authorresponse}

\begin{manuscriptrevision}{\mloc{sec:queries}; 발췌: \quotepages{R3-6}}
본문의 실제 수정 문장은 다음과 같습니다.

\begin{quote}
조건 연결 모델은 committed 위치 Run에서 \texttt{idx<N}이면 \texttt{step()}을 실행하고 인덱스를 늘리는 자기 전이를 하고, \texttt{idx==N}이면 Done으로 이동한다(그림~3).
\end{quote}
\end{manuscriptrevision}

\end{reviewitem}

\clearpage

\begin{reviewitem}{R3-7. 순환 검증과 독립 시험 요구}

\begin{reviewercomment}

6장 합성 시험 설계: 기대한 판정이 나오는 사례만 사전에 선별한 현재 절차는 순환적 검증의 가능성이 있음. 합성 자료 생성과 검사기 개발을 분리하고, 개발에 사용하지 않은 독립 시험 자료를 이용한 평가를 추가.

\end{reviewercomment}

\begin{authorresponse}

사례 선별의 순환성에 대한 지적을 수용했습니다. 기존 40쌍은 기대 결과를 확인한 회귀시험으로 남기고, 144건의 최초 실패 48건 및 개발 후 동일 사례 재시험과 별도로 선별 없는 유한 범위 전수 시험을 추가했습니다.

두 모델 소스·실행 당시 명세·생성기·수식 판정기를 먼저 해시로 고정했습니다. 실행 당시 명세는 해시가 일치하는 재현 입력으로 보존하며 현재 4장은 두 모델의 의미 차이와 적용 범위를 명확히 했습니다. 행동, 의무 생성과 조건 증거의 선언된 조합을 모두 생성했고 기대 판정은 손으로 쓰지 않았습니다. 별도 판정기는 전체 접두열의 최근 증거·생성·해제를 직접 계산하며 상태 갱신 함수를 공유하지 않습니다. 고정 XML에서 추출한 상태 갱신 코드를 C++로 직접 실행한 결과와 수식 판정기의 출력을 상태·누적 문제·최초 사건/의무별로 비교했습니다. 각 모델과 해당 명세의 35,042,928건 대조에서 불일치가 없었습니다. 두 모델 사이의 동등성 시험은 아닙니다. P4 처리 고장의 별도 5,256건에서는 1,884건의 상태 불일치를 모두 보존했습니다. 결과에 따라 사례를 빼거나 정답을 수정하지 않았습니다.

이는 사례 선별 편향과 하나의 계산 경로에 대한 의존을 줄인 부분 대응입니다. 두 경로에 같은 연구 작성자가 관여했고 명세 자체의 공통 오류가 남을 수 있습니다. 유한 상한 밖의 사건열, 독립 자연 의미 정답 및 외부 작성 평가를 입증하지 않으며 UPPAAL 탐색기를 새로 실행한 시험도 아닙니다. 81건 역시 설계자 작성 시험으로 구분합니다. 독립 시험 요구를 완전히 충족했다고 주장하지 않습니다.

\end{authorresponse}

\begin{manuscriptrevision}{\mloc{sec:bounded-results}; 발췌: \quotepages{R3-7}}
본문의 실제 수정 문장은 다음과 같습니다.

\begin{quote}
이 시험은 기대 결과에 맞는 사례만 고르는 절차를 없애고 두 계산 경로를 비교한다. 그러나 수식과 구현이 공유하는 명세의 오류는 함께 남을 수 있다.
\end{quote}
\end{manuscriptrevision}

\end{reviewitem}

\clearpage

\begin{reviewitem}{R3-8. 이력을 보는 비교기}

\begin{reviewercomment}

6장 비교 실험: 현재 사건만 보는 검사기는 정의상 사건 간 오류를 검출할 수 없음. 단순 history-aware rule checker, 유한상태기계 또는 temporal logic 기반 검사 등 보다 의미 있는 비교 방법을 추가할 필요가 있음(옵션).

\end{reviewercomment}

\begin{authorresponse}

완료된 이전 1/3/5개 창과 조건별 FSM 비교를 반영했습니다. 최초 144건에서 일치는 현재 사건 72, 이전 창 84/90/90, 기존 전체 이력 96, FSM144개이며 최초 위반과 간격 0/2/6 결과도 함께 표시했습니다.

현재 사건 검사는 기억 제거 비교입니다. 이전 창은 같은 기존 알고리즘의 창 변형이고 FSM은 조건별 프레임을 받습니다. 단일 의무 투영과 미상 정책의 차이도 결과 원인이므로 단순한 알고리즘 우위로 해석하지 않았습니다. UPPAAL도 FSM보다 정확하다고 주장하지 않습니다.

\end{authorresponse}

\begin{manuscriptrevision}{\mloc{sec:first-results}; 발췌: \quotepages{R3-8}}
본문의 실제 수정 문장은 다음과 같습니다.

\begin{quote}
FSM은 144/144, 최초 위반 72/72였다. UPPAAL이 단순 비교기보다 정확하다는 근거는 없다.
\end{quote}
\end{manuscriptrevision}

\end{reviewitem}

\clearpage

\begin{reviewitem}{R3-9. 선정·교육·합의와 낮은 κ}

\begin{reviewercomment}

6장 표10 및 7.3절 자연 표본 검토: 자동 후보15장면과 비교 장면12개를 선정한 기준, 검토자 자격과 교육, 판정 지침 및 합의 절차를 제시. 특히 Fleiss' κ가0.027로 낮게 나타난 원인을 분석하고, 항목별 판정 분포와 주요 불일치 사례를 함께 제시.

\end{reviewercomment}

\begin{authorresponse}

보존된 원 선정 절차와 재실행으로 재현된 27개 검토 ID를 기술했습니다. 지침의 시간 요구·행동열·텍스트 범주와 외부 근거 미상 처리를 설명하고, 원 CSV의 κ·항목별 일치율/불일치 수·검토자별 텍스트 분포를 추가했습니다. 텍스트는 23/27 불일치, 모순 10/0/1/1·해당 없음 3/5/11/1건으로 라벨 적용 차이가 큽니다.

보존된 지침과 A/B 검토 화면에서 CoC 원문·상대 시각의 제공 범위와 후보군·비교군 구분, 원 장면 ID 및 검사기 결과의 비표시를 확인했습니다. 집단 대응표는 별도 파일에만 있습니다. 이 화면 구성의 확인과 당시 교육·합의 운영 기록의 결손을 구분해 3.3절에 반영했습니다.

원 메모의 시간 요구·조건/순서·선행 관계 정보 부족 등의 의견은 개별 해석으로 제시했습니다. 낮은 κ의 원인을 전문성·교육 탓으로 확정하거나 원인별 건수를 만들지 않았습니다. 자격·교육·중재 및 당시 논의 절차의 기록이 부족하고 합의 사유 메모는 0/27입니다. 따라서 이 자료는 주석 해석의 모호성을 살피는 데 사용하며 독립 자연 모순 정답으로 쓰지 않습니다. 사람 재평가는 미실시입니다.

% Supplemental evidence 2026-10-12
기존 원 라벨을 유지한 AC1 재분석을 추가했습니다(시간 요구/행동열/텍스트 0.627/0.323/0.301). 범주 집합과 원시 일치·기존 κ를 함께 보고, 더 큰 계수만 선택하지 않았습니다. 자격·교육·합의 사유의 기록 결손과 사람 재평가 미실시는 그대로 남깁니다.
% End supplemental evidence

\end{authorresponse}

\begin{manuscriptrevision}{\mloc{sec:human-results}; 발췌: \quotepages{R3-9}}
본문의 실제 수정 문장은 다음과 같습니다.

\begin{quote}
범주 편중에 따른 보정 민감도이며 지침의 안정성이 개선됐다는 뜻이 아니다. 특히 시간 요구는 ``있음''에 판정이 몰려 있어 같은 쌍별 일치율에서도 두 계수의 우연 일치 기대값이 크게 달랐다.
\end{quote}
\end{manuscriptrevision}

\end{reviewitem}

\clearpage

\begin{reviewitem}{R3-10. 40/40 결과와 자연 성능 지표}

\begin{reviewercomment}

7.1절 및 표13: 40/40과20/40 결과는 일반적인 오류 검출 정확도가 아니라 사전에 정의한 규칙의 실행 결과임. 표 제목과 본문에서 이를 ‘검출 성능’으로 오해하지 않도록 표현을 수정하고, 가능하면 실제 모순·비모순 사례를 이용한 precision, recall 및 F1-score를 추가.

\end{reviewercomment}

\begin{authorresponse}

기존표를 ``규칙 실행 회귀시험''으로 바꾸고 P1–P3 주석/구조화 규칙 30쌍과 P4 처리 10쌍으로 분리했습니다. 새 144개의 최초 시험과 보완 재시험도 다른 표로 제시했습니다. 기존 96/144에서 기대 모순 72건 중 미상 출력 18건, 기대 양립 36건 중 미상 출력 30건, 기대 근거 부족 36건 중 미상 출력 36건이라는 혼동 분포를 보고했습니다. 7.2절 첫머리에 기존 UNKNOWN 출력을 근거 부족 범주와 대응시켜 집계하는 규칙을 명시했습니다. 이는 비교를 위한 범주 대응이며, 구문 모호성·복수 의무·대기 중 미상도 누적하는 기존 출력이 해제 근거 부족과 같은 의미라는 뜻은 아님을 함께 밝혔습니다. 양립 사례의 모순 오탐 0/36만으로 성공을 강조하지 않았습니다.

독립적인 실제 모순·비모순 정답은 확보하지 못했으므로 자연 precision, recall, F1은 미산출입니다. 작성 조합 시험의 분류 일치나 최초 위치 일치를 자연 성능으로 바꾸지 않고 모집단 신뢰구간도 만들지 않았습니다.

% Supplemental evidence 2026-10-12
81건의 확정 모순 44/근거 부족 29/양립 8 및 초기화 혼동행렬도 작성 시험으로 보고했습니다. 35,042,928건은 열거한 조합의 수이며 자연 모집단 성능의 분모가 아닙니다. 독립 자연 의미 정답이 없어 자연 P/R/F1을 산출하지 않았습니다.
% End supplemental evidence

\end{authorresponse}

\begin{manuscriptrevision}{\mloc{sec:regression-results}; 발췌: \quotepages{R3-10}}
본문의 실제 수정 문장은 다음과 같습니다.

\begin{quote}
합계 40/40과 20/40을 일반 오류 검출 정확도로 해석하지 않는다.
\end{quote}
\end{manuscriptrevision}

\end{reviewitem}

\clearpage

\begin{reviewitem}{R3-11. 궤적 일치 2·불일치 7·미상 18}

\begin{reviewercomment}

표15의 궤적 불일치 결과: 27개 장면 중 일치2건, 불일치7건, 미상18건으로 나타난 원인을 사례별로 분석. 미상 비율이 높으므로 현재 자료에서 시간적 일관성을 판단할 수 있는 범위도 명확히 기술.

\end{reviewercomment}

\begin{authorresponse}

저장된 사건별 사유와 장면 판정을 재분석한 결과 2/7/18이 재현됐고 장면별 이전 판정과의 차이는 0이었습니다. 본문 \mloc{tab:trajectory-cases}는 판정별 장면 수와 사유별 장면 수(중복 포함)로 정리하고, 27장면 각각의 행동 일치, 종방향 불일치, 문장 해석 미상, 지원 행동 부재, 관측 미상, 단계 시점/해제 근거 부족 사건 건수는 논문 부록의 \mloc{tab:trajectory-details}와 이 답변의 사례표에 수록했습니다. 한 장면에 여러 사유가 겹치며 불일치 사건이 없더라도 미상 사건이 남으면 장면은 미상임을 설명했습니다.

평가 설계 \mloc{sec:aux-design}에 사건별 궤적 판정 규칙과 장면별 집계 우선순위를 추가했습니다. 정지ㆍ감속ㆍ양보는 같은 비진행 행동군으로 비교하고 지원 행동 부재, 문장 해석 실패, 복합 행동의 단계 시점 부족과 관측 행동 불명확은 미상으로 둡니다. 정지ㆍ양보 뒤 가속이 관측돼도 해제 근거가 없으면 미상으로 유지합니다. 장면은 불일치 사건이 있으면 불일치, 불일치 없이 미상 사건이 있으면 미상, 모든 사건이 일치하면 일치로 판정합니다. 궤적의 일치ㆍ불일치ㆍ미상은 4장의 확정 모순ㆍ근거 부족ㆍ양립과 구분했습니다.

조건 해제나 단계 시점의 객체 증거를 속도 변화만으로 확정하지 않았습니다. 불일치 7장면은 문장 간 모순이나 물리 안전의 독립 정답이 아닙니다. 조건 연결 모델에 넣은 작성 증거가 이러한 자연 객체정보 부족을 해결했다는 주장도 하지 않습니다.

% Supplemental evidence 2026-10-12
81건의 작성 경계 결과는 자연 궤적 2/7/18을 대체하지 않습니다. 보관된 장면 연결 후보(13개 구간, 16개 경계)에는 세 조건의 동일 대상·적용 범위·해제 참값 연결이 없어 자연 장면 간 의미 검증도 미실시로 명시했습니다.
% End supplemental evidence

\end{authorresponse}

\begin{manuscriptrevision}{\mloc{sec:trajectory-results}; 발췌: \quotepages{R3-11}}
본문의 실제 수정 문장은 다음과 같습니다.

\begin{quote}
27장면의 궤적 비교는 일치 2, 불일치 7, 미상 18이며 보존된 결과를 다시 실행해도 장면별 판정은 같았다.
\end{quote}
\end{manuscriptrevision}

\clearpage
\textbf{R3-11 상세 사례표 (기존 본문 표 14의 27행 보존)}

익명 G01--G27은 보존 재분석의 행 순서를 따른다. 각 사유는 사건 건수이며 한 장면의 여러 사유 열에 중복될 수 있다. 불일치 사건이 없더라도 미상 사건이 남으면 장면 판정은 미상이다. 이 표의 합은 독립 장면 수나 오류율이 아니다. 본문 \mloc{sec:trajectory-results}와 \mloc{tab:trajectory-cases}의 2/7/18 요약에 대응한다.

\begingroup
\fontsize{8.5pt}{11.5pt}\selectfont
\setlength{\tabcolsep}{4pt}
\begin{longtable}{llrrrrrr}
\caption{R3-11: 27장면별 보존 궤적 판정과 사건 사유}\label{tab:response-trajectory-details}\\
\toprule
장면 & 판정 & \shortstack{행동\\일치} & \shortstack{종방향\\불일치} & \shortstack{해석\\미상} & \shortstack{행동\\미지원} & \shortstack{관측\\미상} & \shortstack{시점/해제\\부족} \\
\midrule
\endfirsthead
\multicolumn{8}{l}{R3-11 사례별 사유 (계속)}\\
\toprule
장면 & 판정 & \shortstack{행동\\일치} & \shortstack{종방향\\불일치} & \shortstack{해석\\미상} & \shortstack{행동\\미지원} & \shortstack{관측\\미상} & \shortstack{시점/해제\\부족} \\
\midrule
\endhead
\bottomrule
\endfoot
G01 & 미상 & 0 & 0 & 0 & 4 & 2 & 0 \\
G02 & 미상 & 0 & 0 & 2 & 1 & 0 & 0 \\
G03 & 불일치 & 1 & 1 & 2 & 0 & 0 & 1 \\
G04 & 미상 & 1 & 0 & 1 & 1 & 0 & 1 \\
G05 & 불일치 & 0 & 1 & 1 & 4 & 0 & 0 \\
G06 & 미상 & 2 & 0 & 0 & 0 & 0 & 2 \\
G07 & 미상 & 3 & 0 & 1 & 0 & 0 & 0 \\
G08 & 미상 & 4 & 0 & 0 & 3 & 0 & 2 \\
G09 & 미상 & 3 & 0 & 0 & 1 & 0 & 2 \\
G10 & 불일치 & 0 & 1 & 0 & 1 & 0 & 0 \\
G11 & 미상 & 3 & 0 & 1 & 1 & 0 & 0 \\
G12 & 미상 & 3 & 0 & 1 & 0 & 0 & 2 \\
G13 & 불일치 & 1 & 1 & 1 & 0 & 0 & 1 \\
G14 & 미상 & 1 & 0 & 0 & 0 & 0 & 2 \\
G15 & 미상 & 0 & 0 & 0 & 4 & 0 & 1 \\
G16 & 일치 & 3 & 0 & 0 & 0 & 0 & 0 \\
G17 & 미상 & 0 & 0 & 0 & 2 & 0 & 1 \\
G18 & 일치 & 3 & 0 & 0 & 0 & 0 & 0 \\
G19 & 미상 & 0 & 0 & 0 & 2 & 0 & 1 \\
G20 & 미상 & 0 & 0 & 4 & 3 & 0 & 0 \\
G21 & 미상 & 0 & 0 & 0 & 2 & 0 & 0 \\
G22 & 미상 & 1 & 0 & 0 & 3 & 0 & 0 \\
G23 & 불일치 & 4 & 1 & 0 & 0 & 0 & 0 \\
G24 & 불일치 & 2 & 3 & 0 & 0 & 0 & 0 \\
G25 & 불일치 & 0 & 1 & 1 & 1 & 0 & 1 \\
G26 & 미상 & 0 & 0 & 0 & 6 & 0 & 0 \\
G27 & 미상 & 2 & 0 & 0 & 2 & 0 & 0 \\
\end{longtable}
\endgroup


일치 G16/G18은 지원 행동과 기록 궤적이 맞는 사례다. 불일치 G03, G05, G10, G13, G23, G24, G25에는 종방향 행동 불일치가 있다. 나머지 18장면은 위 표의 해석ㆍ지원 행동ㆍ관측ㆍ시점/해제 근거가 충분하지 않다. 예를 들어 G06은 행동 일치 2건이 있어도 시점/해제 부족 2건 때문에 미상이고, G01은 행동 미지원 4건과 관측 미상 2건이 남는다. 속도 정보만으로 객체의 해제 상태를 채우지 않았으며 불일치를 문장 모순이나 안전 위반 정답으로 사용하지 않는다.

\end{reviewitem}

\section*{맺음말}

심사위원님의 의견에 따라 관련 연구의 비교, 자료 선정 절차, P4 처리 오류의 분리, 미상과 장면 경계의 처리, 입력 변환 사례, UPPAAL 상태와 질의, 이력을 보는 비교기, 검토 절차와 궤적 판정 기준을 보완했습니다. 순환 검증의 우려에 대해서는 선별 없는 유한 범위 전수 대조를 추가했으나, 개발에 쓰지 않은 외부 독립 시험과 자연 오류 정답은 확보하지 못했으므로 이를 본문의 적용 범위와 후속 과제로 명시했습니다.

\end{document}
